% ===== 原文 pp.262–276（PDF p.277–291）=====
% 第 6 章：§6.2（§6.2.1 末 → §6.2.5 开头）
% 章标题与 §6.2 标题由 c06a 写，本块不写。
% 块首续 c06a 末尾 Lemma 6.9 的证明：普通段落续写，末尾补 \hfill$\square$，不重开 proof 环境。
% 本块无编号公式（tag 区间为空）。定理类环境 16 个（6.10–6.25）；图 Figure 5–11（p.265–p.276），已裁 work/figures/f6_5.png … f6_11.png。
% 依赖：preamble 已加 \newtheorem{assumption}[theorem]{假设}（Lead 已确认）。

% [原文 p.262]
为证明这一无关性，首先注意：若 $\tau\subseteq\sigma\in\mathcal{P}_{\max}$，$\tau'\subseteq\sigma'\in\mathcal{P}_{\max}$，则 $I^k_{\omega,\sigma}\subseteq I^k_{\omega,\tau}$、$I^k_{\omega,\sigma'}\subseteq I^k_{\omega,\tau'}$，故 $c_{p,\sigma}=c_{p,\tau}$、$c_{p,\sigma'}=c_{p,\tau'}$。因此只需证明：当 $p\notin I^k_{\omega,\sigma}$、$I^k_{\omega,\sigma'}$ 时 $c_{p,\sigma}=c_{p,\sigma'}$。
首先考虑存在余维一胞腔 $\rho\in\mathcal{P}$ 使 $\sigma\cap\sigma'=\rho$ 的情形。则 $I^k_{\omega,\rho}=I^k_{\omega,\sigma}\cup I^k_{\omega,\sigma'}$，故 $p\notin I^k_{\omega,\rho}$。于是 $c_{p,\sigma}=c_{p,\rho}=c_{p,\sigma'}$。
对一般情形，考虑函数 $\varphi_\omega$ 的 Newton 多面体 $\Delta_{\varphi_\omega}\subseteq M_\R$（见定义 1.8）。则 $\Delta_{\varphi_\omega}$ 有顶点 $\check\sigma$、$\check\sigma'$，其中 $\check\sigma=-m_\sigma$、$\check\sigma'=-m_{\sigma'}$。设 $p=(n,r)$。于是 $p\notin I^k_{\omega,\sigma}\cup I^k_{\omega,\sigma'}$ 这一事实说明
\[
r+\langle-m_\sigma,n\rangle<k+1,\qquad r+\langle-m_{\sigma'},n\rangle<k+1,
\]
即 $\check\sigma$ 与 $\check\sigma'$ 都位于 $M_\R$ 中由 $\langle\cdot,n\rangle=k+1-r$ 定义的超平面的同一侧。
由于 $\Delta_{\varphi_\omega}$ 是凸的，可以找到一列顶点 $\check\sigma=\check\sigma_1,\check\sigma_2,\dots,\check\sigma_n=\check\sigma'$，它们都在该超平面的同一侧，并且 $\check\sigma_{i-1}$ 与 $\check\sigma_i$ 由 $\Delta_{\varphi_\omega}$ 的一条边 $\check\rho_i$ 相连，而 $\check\rho_i$ 对应于一个含 $\omega$ 的余维一胞腔 $\rho_i$。则对 $1\le i\le n$ 有 $p\notin I^k_{\omega,\sigma_i}$，而 $\sigma_{i-1}\cap\sigma_i=\rho_i$。于是我们处于已考虑过的第一种情形，得到 $c_{p,\sigma}=c_{p,\sigma_2}=\cdots=c_{p,\sigma'}$，正如所欲。
\hfill$\square$
% [Tutor Note: 本块开头（p.262）是 c06a 末尾 Lemma 6.9 证明的续写；按 Lead 裁决以普通段落续写并以 $\hfill\square$ 封口，未重开 proof 环境。核对锚点：符号 $\check\sigma=-m_\sigma$、$\check\sigma'=-m_{\sigma'}$ 与超平面 $\langle\cdot,n\rangle=k+1-r$（PDF p.277 图像）。]

由于
\[
U^k_\omega=\operatorname{Spec}R_\omega/(t^{k+1}),
\]
我们看到开集 $U^k_\omega$ 可以从环 $R^k_{\omega,\tau}$ 重构出来，而这些环是 $\pi^{-1}(0)$ 的分层的仿射子集的加厚的坐标环，这里 $\pi:\mathbb{P}_{\widetilde\Delta}\to\mathbb{A}^1_\Bbbk$ 是通常的投影。当然，$X_k$ 于是可以从开集 $U^k_\omega$ 重构出来，因为若 $\omega_1\subseteq\omega_2$，我们有自然包含 $U^k_{\omega_2}\subseteq U^k_{\omega_1}$。这些包含由环的自然包含
\[
R_{\omega_1}/(t^{k+1})\hookrightarrow R_{\omega_2}/(t^{k+1})
\]
诱导，而后者由包含 $P_{\omega_1}\subseteq P_{\omega_2}$ 诱导。用逆向系统的语言，这可以看作由自然包含 $R^k_{\omega_1,\tau}\subseteq R^k_{\omega_2,\tau}$（只要 $\omega_1\subseteq\omega_2\subseteq\tau$）诱导。无论如何，通过这些自然包含把开集 $U^k_\omega$ 粘合起来，就得到概形 $X_k$。
要点在于：虽然 $X_k$ 的这个构造看起来更复杂，但我们现在给出了 $X_k$ 的一个局部得多的描述，它可用于更一般的设定，这正是我们接下来要探讨的。

\subsection{整体情形：无奇点}
推广上述方法的第一步是把 $B$ 换成一个一般的、无奇点的整热带流形（我们仍然允许 $B$ 有边界或非紧）。于是考虑这样的 $B$，或者更确切地说，考虑三元组 $(B,\mathcal{P},\varphi)$。注意现在我们允许 $\varphi$ 是多值的。这很重要：第一，我们实际上只需要 $\varphi$ 的局部描述（精确到线性函数）；第二，一般不存在整体的严格凸单值分片线性函数。

\begin{example}
设 $\Gamma\subseteq N$ 是子格，并令 $B=N_\R/\Gamma$。我们取一个多面体分解 $\mathcal{P}$ 与一个具有整斜率的多值严格凸分片线性函数。在此情形，这些数据可以放在 $B$ 的泛覆叠 $N_\R$ 上看：此时 $\mathcal{P}$ 应看作 $N_\R$ 的一个 $\Gamma$-周期多面体分解，而 $\varphi:N_\R\to\R$ 应是严格凸的，并对任意 $\gamma\in\Gamma$ 满足
\[
\varphi(x+\gamma)=\varphi(x)+\alpha_\gamma(x)
\]
其中 $\alpha_\gamma:N_\R\to\R$ 是某个仿射线性函数。
\end{example}

% [原文 p.263]
事实上，我们即将说的内容即使对带奇点的 $B$ 也适用，所以眼下我们考虑最一般的情形：$(B,\mathcal{P},\varphi)$，其中 $B$ 是带奇点的整热带流形。

\begin{definition}
给定如上 $(B,\mathcal{P},\varphi)$。设 $\{(U_i,\varphi_i)\}$ 是 $\varphi$ 在 $B$ 的开覆盖 $\{U_i\}$ 上的一组代表。按如下方式把 $B_0$ 上的层 $\mathcal{P}_\varphi$ 定义为扩张
\[
0\to\Z\to\mathcal{P}_\varphi\to\Lambda\to0 :
\]
首先，$\mathcal{P}_\varphi|_{U_i}\cong\Z\oplus\Lambda|_{U_i}$。在 $x\in U_i\cap U_j\cap B_0$ 处，我们把 $(r,m)\in(\Z\oplus\Lambda|_{U_i})_x$ 与 $(r+d(\varphi_j-\varphi_i)(m),m)\in(\Z\oplus\Lambda|_{U_j})_x$ 等同，从而给出 $\Z\oplus\Lambda|_{U_i}$ 与 $\Z\oplus\Lambda|_{U_j}$ 在 $U_i\cap U_j\cap B_0$ 上的粘合。
注意这是合理的：$\varphi_j-\varphi_i$ 是具整斜率的仿射线性函数，故 $d(\varphi_j-\varphi_i)$ 自然地是 $U_i\cap U_j\cap B_0$ 上 $\check\Lambda$ 的一个截面，于是可以在切向量 $m$ 上取值。
对 $\mathcal{P}_\varphi$ 的截面 $m$，我们把它在投影 $\mathcal{P}_\varphi\to\Lambda$ 下的像记作 $\bar m$。
\end{definition}

\begin{remark}
事实上，层 $\mathcal{P}_\varphi$ 可以描述为 $\operatorname{Aff}(\check B,\Z)$，即 $\check B$（$B$ 的离散 Legendre 变换）上仿射线性函数的层，只要把 $B$ 与 $\check B$ 作一个典范等同。这一描述在 [49] 中被使用；这里我们用上述描述，以免过多使用离散 Legendre 变换。
\end{remark}

我们接着把幺半群 $P_\tau$ 的概念推广到这一情形。此前它们生活在 $N\oplus\Z$ 中。把 $N$ 看作 $N_\R$ 的整切向量空间，则自然地可以想象 $N\oplus\Z$ 的整体版本应是 $\Lambda$ 的某个由 $\Z$ 扩张而来的扩张。事实上，$\mathcal{P}_\varphi$ 正是正确的扩张。

\begin{definition}
点 $x\in B_0$ 处的一个指数（exponent）是 $\mathcal{P}_\varphi$ 在 $x$ 处的茎的元素。对任意点 $x\in B_0$，我们定义幺半群
\[
\mathcal{P}_{\varphi,x}\subseteq\mathcal{P}_{\varphi,x},
\]
使 $\mathcal{P}^{gp}_{\varphi,x}=\mathcal{P}_{\varphi,x}$，如下。设 $x\in U_i$。对每个满足 $x\in\sigma$ 的 $\sigma\in\mathcal{P}_{\max}$，设 $\varphi_{i,\sigma}\in\check\Lambda_x$ 是 $\varphi_i|_\sigma$ 的微分。设 $\Sigma_x$ 是 $\Lambda_x\otimes_\Z\R$ 中由切楔
\[
\Sigma_x:=\{T_x\tau\mid x\in\tau\in\mathcal{P}\}
\]
给出的扇（其锥不必严格凸）。设 $\varphi_{i,x}:|\Sigma_x|\to\R$ 是在锥 $T_x\sigma$（$\sigma\in\mathcal{P}_{\max}$）上由 $\varphi_{i,\sigma}$ 给出的函数。定义
\[
\mathcal{P}_{\varphi,x}:=\{(r,m)\mid m\in|\Sigma_x|,\ r\ge\varphi_{i,x}(m)\}\subseteq\mathcal{P}_{\varphi,x},
\]
这里用了分裂 $\mathcal{P}_{\varphi,x}=\Z\oplus\Lambda|_{U_i}$。
\end{definition}

% [Tutor Note: 定义 6.11（p.263）的粘合公式为：在 $x\in U_i\cap U_j\cap B_0$ 处把 $(r,m)$ 认同为 $(r+d(\varphi_j-\varphi_i)(m),m)$；这里 $\bar m$ 表示 $m$ 在投影 $\mathcal{P}_\varphi\to\Lambda$ 下的像。注意本书中多面体分解 $\mathcal{P}$ 与层 $\mathcal{P}_\varphi$ 用同一个字母（花体 P），仅靠下标 $\varphi$ 区分，转录时勿混。]
% [Tutor Note: 定义 6.13 中“$\mathcal{P}_{\varphi,x}\subseteq\mathcal{P}_{\varphi,x}$”看似重复，实为原文写法：左边是 $x$ 处的幺半群、右边是层 $\mathcal{P}_\varphi$ 在 $x$ 处的茎；同理 $\mathcal{P}^{gp}_{\varphi,x}=\mathcal{P}_{\varphi,x}$ 中左边表示该幺半群的群化。另 $\Sigma_x$ 是“其锥不必严格凸”的扇，见定义中括号内的说明。]
% [原文 p.264]
层 $\mathcal{P}_\varphi$ 的要点在于：幺半群 $\mathcal{P}_{\varphi,x}$ 的定义不依赖于用来定义它的开集 $U_i$ 的选择。确实，注意
\[
\varphi_{j,x}=\varphi_{i,x}+d(\varphi_j-\varphi_i).
\]
于是若 $m\in|\Sigma_x|$，有
\begin{align*}
(r,m)\in\mathcal{P}_{\varphi,x}\ \text{（在 }U_i\text{ 中定义）}
&\iff r\ge\varphi_{i,x}(m)\\
&\iff r\ge\varphi_{j,x}(m)-d(\varphi_j-\varphi_i)(m)\\
&\iff r+d(\varphi_j-\varphi_i)(m)\ge\varphi_{j,x}(m)\\
&\iff (r+d(\varphi_j-\varphi_i)(m),m)\in\mathcal{P}_{\varphi,x}\ \text{（在 }U_j\text{ 中定义）}.
\end{align*}
这表明 $\mathcal{P}_{\varphi,x}$ 是良定义的。此外，同样的论证表明定义 6.7 中定义的阶的概念也是有意义的。若我们使用如上由 $\varphi_i$ 定义的分裂，并把 $p\in\mathcal{P}_{\varphi,x}$ 在该分裂下写成 $(r,m)$，则对满足 $x\in\sigma$ 的 $\sigma\in\mathcal{P}_{\max}$，定义
\[
\operatorname{ord}_\sigma(p)=r-\varphi_{i,\sigma}(m),
\]
容易如上验证它不依赖于分裂的选择。
类似地，对 $\sigma\in\mathcal{P}^\partial_{\max}$（即 $\mathcal{P}$ 中含于 $\partial B$、维数为 $\dim B-1$ 的胞腔的集合），若 $x\in\sigma$，我们定义
\[
\operatorname{ord}_\sigma(p)=\langle n_\sigma,m\rangle,
\]
其中 $n_\sigma\in\check\Lambda_x$ 是本原的、在 $T_x\sigma$ 上消失、且在 $|\Sigma_x|$ 上非负。这显然不依赖于分裂，只依赖于 $m=\bar p$。

\begin{remark}
虽然本小节我们将集中于 $B$ 无奇点的情形，但让我们描述当 $\dim B=2$ 且 $B$ 有正奇点时 $\mathcal{P}$ 在一个奇点附近的局部单值性。设 $p\in B$ 是这样一个奇点，$p\in\omega$ 是一维面。则用 $U_\omega$（$\omega$ 的开星）上的代表 $\varphi_{U_\omega}$，得到分裂
\[
\mathcal{P}_\varphi|_{U_\omega}=\Z\oplus\Lambda|_{U_\omega}.
\]
于是 $\mathcal{P}_\varphi$ 绕图 1 中所描绘的环路的单值性分裂为
\[
(r,\bar m)\mapsto\bigl(r,\bar m+\kappa_\omega\langle\check d_\omega,\bar m\rangle d_\omega\bigr)
=(r,\bar m)+\kappa_\omega\langle\check d_\omega,\bar m\rangle(0,d_\omega).
\]
现在回忆：多值分片线性函数的代表 $\varphi_{U_\omega}$ 由 $\varphi_{U_\omega}=\lambda+\varphi_\omega\circ S_\omega$ 给出（定义 1.30），其中 $\lambda$ 在 $U_\omega$ 上是仿射线性的。然而微分 $d\lambda$ 是 $\check\Lambda$ 的截面。由 $\Lambda$ 的单值性描述容易看出，转置单值性在 $\check\Lambda$ 上的作用为
\[
n\mapsto n+\kappa_\omega\langle n,d_\omega\rangle\check d_\omega.
\]
因此特别地，若 $\lambda$ 在 $U_\omega$ 上良定义，则 $d\lambda$ 必须是单值不变的。所以若 $\kappa_\omega\neq0$，就必须有 $\langle d\lambda,d_\omega\rangle=0$。于是特别地 $\varphi_{U_\omega}$ 在 $\omega$ 上是常数，而对 $x\in\operatorname{Int}(\omega)\setminus\{p\}$，$\varphi_{U_\omega,x}$ 沿 $\omega$ 的切空间（由 $d_\omega$ 生成）为零。于是 $(0,d_\omega)\in\Z\oplus\Lambda_{v_+}=\mathcal{P}_{\varphi,v_+}$ 事实上可以解释为 $\mathcal{P}_{\varphi,v_+}$ 的唯一元素

% [原文 p.265]
\begin{figure}[htbp]
\centering
\includegraphics[width=0.42\textwidth]{f6_5.png}
\figtitle{5}{}
\end{figure}
$m_-$，使得 $\bar m_-=d_\omega$ 且 $\operatorname{ord}_{\sigma_\pm}(m_-)=0$。于是 $\mathcal{P}_{\varphi,v_+}$ 上的单值性为
\[
m\mapsto m+\kappa_\omega\langle\check d_\omega,\bar m\rangle m_-,
\]
而这个描述现在不依赖于分裂的选择。\hfill$\square$
\end{remark}

% [Tutor Note: 注 6.14（p.264–265）中的单值性公式为 $(r,\bar m)\mapsto(r,\bar m)+\kappa_\omega\langle\check d_\omega,\bar m\rangle(0,d_\omega)$，其简化形式 $m\mapsto m+\kappa_\omega\langle\check d_\omega,\bar m\rangle m_-$ 出现在 p.265；$m_-$ 由 $\bar m_-=d_\omega$ 与 $\operatorname{ord}_{\sigma_\pm}(m_-)=0$ 唯一确定。文中的“图 1”指第 6 章 Figure 1（在 c06a 页段），不是本块图号。]
现在，对 $\omega\subseteq\tau\subseteq\sigma$（其中 $\omega,\tau\in\mathcal{P}$，$\sigma\in\mathcal{P}_{\max}$），我们将如下定义环 $R^k_{\omega,\tau,\sigma}$。取任意点 $x\in\operatorname{Int}(\omega)\cap B_0$；这给出幺半群 $\mathcal{P}_{\varphi,x}\subseteq\mathcal{P}_{\varphi,x}$。我们取一条从 $x$ 出发、立即进入 $\sigma$ 的内部、终点为 $y$ 的短路，见图 5。通过局部系统 $\mathcal{P}_\varphi$ 中的平行移动，我们把 $\mathcal{P}_{\varphi,x}$ 与 $\mathcal{P}_{\varphi,y}$ 等同。幺半群 $\mathcal{P}_{\varphi,x}\subseteq\mathcal{P}_{\varphi,x}$ 于是通过这一等同给出一个包含在 $\mathcal{P}_{\varphi,y}$ 中的幺半群，我们称之为 $\mathcal{P}_{\varphi,\omega,\sigma}$。容易看出它与 $x$ 的选择无关，因为层 $\mathcal{P}_\varphi$ 在可缩集 $\sigma\cap B_0$ 上没有单值性。
这里的基本想法是我们用 $\sigma$ 而不是 $x$ 作为参照；这在后面的记账中会更好。
现在我们定义一个理想
\[
I^k_{\omega,\tau,\sigma}\subseteq\mathcal{P}_{\varphi,x}=\mathcal{P}_{\varphi,\omega,\sigma},
\]
办法是
\[
I^k_{\omega,\tau,\sigma}=\{p\in\mathcal{P}_{\varphi,x}\mid \exists\sigma'\in\mathcal{P}_{\max}\cup\mathcal{P}^\partial_{\max}\ \text{满足}\ \tau\subseteq\sigma'\ \text{s.t.}\ \operatorname{ord}_{\sigma'}(p)>k\}.
\]
容易验证它在 $\mathcal{P}_{\varphi,\omega,\sigma}$ 中是良定义的，不依赖于 $x$ 的选择。我们于是定义
\[
R^k_{\omega,\tau,\sigma}:=\Bbbk[\mathcal{P}_{\varphi,\omega,\sigma}]/I^k_{\omega,\tau,\sigma}.
\]
这与上一小节的 Mumford 退化完全类似。
正如我们所说，到目前为止，即使 $B$ 有奇点这也成立。然而本小节的剩余部分，我们需要假设 $B$ 没有奇点。
在此情形，给定 $\omega\subseteq\tau\subseteq\sigma$、$\sigma'\in\mathcal{P}_{\max}$，存在典范同构
\[
R^k_{\omega,\tau,\sigma}\cong R^k_{\omega,\tau,\sigma'}.
\]
它由沿连接 $y\in\sigma$ 与 $y'\in\sigma'$ 的路径（如图 6）在 $\mathcal{P}_\varphi$ 中的平行移动得到。于是这些环不依赖于参照胞腔，我们就照 Mumford 情形那样把它们写成 $R^k_{\omega,\tau}$。

% [原文 p.266]
\begin{figure}[htbp]
\centering
\includegraphics[width=0.42\textwidth]{f6_6.png}
\figtitle{6}{比较以不同参照胞腔给出的环}
\end{figure}

我们于是构造开集
\[
U^k_\omega:=\operatorname{Spec}\varprojlim R^k_{\omega,\tau},
\]
这里对 $\omega\subseteq\tau_1\subseteq\tau_2$ 用了典范满射
\[
\psi_{\tau_1,\tau_2}:R^k_{\omega,\tau_2}\to R^k_{\omega,\tau_1}.
\]
此外，对 $\omega_1\subseteq\omega_2$，我们有包含
\[
U^k_{\omega_2}\subseteq U^k_{\omega_1},
\]
它们由包含 $R^k_{\omega_1,\tau}\subseteq R^k_{\omega_2,\tau}$ 诱导，于是我们可以用这些等同把各开集 $U^k_\omega$ 粘合起来。
这给出 $\mathcal{O}_k$ 上的一个概形 $X_k$。若取 $k=0$，则不难验证 $X_0=\check X_0(B,\mathcal{P})$，即 $X_0$ 由把 $\mathbb{P}_\sigma$（$\sigma\in\mathcal{P}_{\max}$）粘合而成。确实，考虑 $\sigma\in\mathcal{P}_{\max}$。对 $\omega\subseteq\sigma$ 为面，设 $x\in\operatorname{Int}(\omega)$ 是一点，并设 $x\in U_i$（$\varphi$ 由 $\varphi_i$ 代表），于是
\[
\mathcal{P}_{\varphi,x}=\{(r,m)\mid m\in|\Sigma_x|,\ r\ge\varphi_{i,x}(m)\}.
\]
然后把它与 $\mathcal{P}_{\varphi,\omega,\sigma}$ 等同，我们看到
\[
\mathcal{P}_{\varphi,\omega,\sigma}\setminus I^0_{\omega,\sigma,\sigma}
=\{(r,m)\mid m\in|\Sigma_x|,\ r\ge\varphi_{i,x}(m),\ r=\varphi_{i,\sigma}(m)\}.
\]
于是 $(r,m)\in\mathcal{P}_{\varphi,\omega,\sigma}\setminus I^0_{\omega,\sigma,\sigma}$ 当且仅当 $m$ 指向 $\sigma$ 的内部且 $r=\varphi_{i,\sigma}(m)$，这由 $\varphi_{i,x}$ 的严格凸性得到。所以我们可以把 $\Bbbk[\mathcal{P}_{\varphi,\omega,\sigma}]/I^0_{\omega,\sigma,\sigma}$ 与 $\Bbbk[(T_\omega\sigma)\cap\Lambda_x]$ 等同。把 $\Lambda_x$ 与 $\Lambda_y$（$y\in\operatorname{Int}(\sigma)$ 为固定点）等同，我们看到 $\operatorname{Spec}\Bbbk[(T_\omega\sigma)\cap\Lambda_y]$ 是 $\mathbb{P}_\sigma$ 中由锥 $N_{\sigma\omega}\in\check\Sigma_\sigma$（$\sigma$ 的法扇）所指定的仿射开子集。此外，当 $\omega_1\subseteq\omega_2\subseteq\sigma$ 时的包含 $R^0_{\omega_1,\sigma}\subseteq R^0_{\omega_2,\sigma}$ 对应于自然包含
\[
\Bbbk[(T_{\omega_1}\sigma)\cap\Lambda_y]\subseteq\Bbbk[(T_{\omega_2}\sigma)\cap\Lambda_y],
\]
这说明当我们通过这些包含把集合 $\operatorname{Spec}R^0_{\omega,\sigma}$ 粘合起来时，得到 $\mathbb{P}_\sigma$。由于集合 $\operatorname{Spec}R^0_{\omega,\sigma}$ 是 $U^0_\omega$ 的不可约分支，我们看到 $X_0$ 确实是环面簇 $\mathbb{P}_\sigma$ 的粘合。
所以这表明：目前的构造已经足够局部，可以处理任意的无奇点整热带流形。

% [原文 p.267]
\begin{figure}[htbp]
\centering
\includegraphics[width=0.40\textwidth]{f6_7.png}
\figtitle{7}{多面体 $\widetilde\Delta$}
\end{figure}

\begin{example}
设 $B=\R/d\Z$，其中 $d$ 为正整数，并取
\[
\mathcal{P}=\{[i,i+1]\mid 1\le i\le d\}\cup\bigl\{\{i\}\mid 1\le i\le d\bigr\}.
\]
我们可以在 $B$ 的泛覆叠 $\R$ 上整体地定义 $\varphi$ 为
\[
\varphi(x)=ix-\frac{i(i+1)}{2}\quad(x\in[i,i+1]),
\]
于是 $\varphi$ 满足周期性条件
\[
\varphi(x+d)=\varphi(x)+d\cdot x+\frac{d(d-1)}{2}.
\]
在此情形，$X_k$ 是由 $d$ 条有理曲线组成的一个循环的 $k$ 阶光滑化。事实上，我们并不需要使用这个局部描述。有一个在 $B$ 的泛覆叠上工作的 $X_k$ 的整体构造，如下。
令
\[
\widetilde\Delta:=\{(n,r)\in\R\oplus\R\mid r\ge\varphi(n)\},
\]
如图 7 所示。设 $\check\Sigma_{\widetilde\Delta}$ 是 $\widetilde\Delta$ 的法扇。这是一个无限扇，其射线由 $(i,1)$（$i\in\Z$）生成，二维锥由 $\{(i,1),(i+1,1)\}$（$i\in\Z$）生成，见图 8。
\begin{figure}[htbp]
\centering
\includegraphics[width=0.62\textwidth]{f6_8.png}
\figtitle{8}{$\widetilde\Delta$ 的法扇}
\end{figure}
% [Tutor Note: 图 7 与图 8 均裁自 PDF p.282（原书 p.267）：图 7 是 $\widetilde\Delta=\{(n,r)\in\R\oplus\R\mid r\ge\varphi(n)\}$ 的图形（灰色区域），图 8 是它的法扇 $\check\Sigma_{\widetilde\Delta}$（射线 $(i,1)$）。两图在本块均无公式依赖，按原图裁剪。]
这定义了环面簇 $X:=X_{\check\Sigma_{\widetilde\Delta}}$，它不是有限型的，连同由投影到第二坐标给出的正则函数 $\pi:X\to\mathbb{A}^1_\Bbbk$。这只是 Mumford 构造的无限版本：$\pi^{-1}(0)$ 是 $\mathbb{P}^1$ 的无限链，而 $\pi$ 的一般纤维是 $\mathbb{G}_m$ 的一个拷贝。注意群 $\Z$ 线性地作用在扇 $\check\Sigma_{\widetilde\Delta}$ 上，由
\[
\Z\ni1\longmapsto\begin{pmatrix}1&d\\0&1\end{pmatrix}\in\mathrm{GL}_2(\Z).
\]

% [原文 p.268]
于是 $\Z$ 作用在 $X$ 上，而这个作用限制到 $\pi^{-1}(0)$ 上时，是把 $\mathbb{P}^1$ 的链平移 $d$ 个位置。
理想情况下，人们希望除以这个作用，定义 $X/\Z$。然而这个商并不是好行为的。确实，可以这样看出：若取 $t$ 为 $\mathbb{A}^1$ 上的坐标，并取 $X$ 的大环面轨道上对应于 $\R\oplus\R$ 的标准基的坐标 $z_1,z_2$，则 $\pi^{-1}(t)$ 被认同为具坐标 $z_1$ 的 $\mathbb{G}_m$，且 $z_2=t$。$1\in\Z$ 的作用在这些坐标下是 $z_1\mapsto z_1z_2^d$、$z_2\mapsto z_2$，从而在 $\mathbb{G}_m=\pi^{-1}(t)$ 上的作用是 $z_1\mapsto z_1t^d$。当 $t=1$ 时这个作用有不动点。因此这个商不是好行为的，甚至在复流形的范畴中都不存在。
有两种办法修正这一点。较容易的办法是在复解析范畴中工作，此时注意 $\pi^{-1}(D)/\Z$ 是有意义的，其中 $D=\{t\mid|t|<1\}$ 是单位圆盘。这给出一个纤维化 $\pi^{-1}(D)/\Z$，其一般纤维是同构于 $\C/\langle1,\frac{1}{2\pi\ii}\log t\rangle$ 的椭圆曲线，而其在 $0$ 上的纤维是 $d$ 条有理曲线的循环。注意这个纤维化不能延拓为 $\mathbb{A}^1$ 上的纤维化。
另一种办法是令
\[
\widetilde X_k:=X\times_{\mathbb{A}^1}\mathcal{O}_k,
\]
然后 $X_k:=\widetilde X_k/\Z$ 是有意义的。事实上，$X_k$ 恰好就是我们上面构造的 $\check X_0(B,\mathcal{P})$ 的 $k$ 阶加厚。取 $k\to\infty$ 的极限，就得到一个定义在 $\Bbbk[[t]]$ 的形式谱上的形式概形。适当小心地处理后，可以应用 Grothendieck 存在性定理（[52], 5.4.5）证明：它来自某个定义在 $\operatorname{Spec}\Bbbk[[t]]$ 上的概形沿中心纤维的完备化。在代数几何设定中，这已是我们一般所能达到的最好结果。
更一般地，若 $B=\R^g/\Gamma$，其中 $\Gamma\subseteq\Z^g\subseteq\R^g$ 是格，同样的程序也有效。这即 Mumford 的阿贝尔簇退化构造，见 [84] 与 [2]；刚刚讨论的一维情形通常称为 Tate 曲线。所以事实上，在引入奇点之前，这个构造并没有真正给我们带来新东西。\hfill$\square$
\end{example}

\subsection{引入奇点：策略}
现在考虑允许 $B$ 有奇点的情形。此时我们限于 $\dim B=2$ 的情形，因为这是我们要覆盖的情形。[49] 的许多复杂性来自高维的困难。在二维情形，[49] 的方法可以按照与 Kontsevich 与 Soibelman 在 [70] 中的论证大致相同的方式来描述。
于是考虑一个局部情形：我们有一条边 $\omega\in\mathcal{P}$ 含有一个奇点，并假定它是正的且简单的。按图 9 的标号，以点 $x$ 为基点、顺时针绕过奇点的环路 $\gamma$ 的单值性取形式
\[
T_\gamma(m)=m+\langle\check d_\omega,m\rangle d_\omega,
\]
其中 $d_\omega$ 是从 $v_+$ 指向 $v_-$ 的本原切向量，而 $\check d_\omega$ 是本原的、与 $d_\omega$ 正交、且在 $\sigma_+$ 上为正。于是如果我们取 $\Lambda_x$ 的一组基 $e_1=d_\omega$ 与 $e_2$（指向 $\sigma_+$ 内部的本原切向量），就得到
\[
T_\gamma(e_1)=e_1,\qquad T_\gamma(e_2)=e_1+e_2.
\]

% [原文 p.269]
\begin{figure}[htbp]
\centering
\includegraphics[width=0.46\textwidth]{f6_9.png}
\figtitle{9}{}
\end{figure}

此外，我们可以取分片线性函数 $\varphi$ 在 $\sigma_-$ 上为 $0$，在 $\sigma_+$ 上斜率为 $e_2^*$，其中 $e_1^*,e_2^*$ 是 $\check\Lambda_x$ 的对偶基。
现在我们有了环 $R^k_{\omega,\sigma_\pm,\sigma_\pm}$ 与 $R^k_{\omega,\omega,\sigma_\pm}$。按图 9 所指出的从 $y_\pm$ 到 $x$ 的路径把 $\Lambda_{y_\pm}$ 与 $\Lambda_x$ 等同，则记 $\Z^3=\Z\oplus\Lambda_x=\mathcal{P}_{\varphi,x}$，可把 $\mathcal{P}_{\varphi,x}$ 的生成元写成 $(1,0,1)$、$(0,0,-1)$ 与 $(0,\pm1,0)$。令
\[
u=z^{(1,0,1)},\qquad v=z^{(0,0,-1)},\qquad w=z^{(0,1,0)},
\]
于是
\begin{align*}
R^k_{\omega,\sigma_+,\sigma_+}&\cong\Bbbk[u,v,w^{\pm1}]/(v^{k+1}),\\
R^k_{\omega,\sigma_-,\sigma_-}&\cong\Bbbk[u,v,w^{\pm1}]/(u^{k+1}),\\
R^k_{\omega,\omega,\sigma_\pm}&\cong\Bbbk[u,v,w^{\pm1}]/(u^{k+1},v^{k+1}).
\end{align*}
现在我们有典范满射
\[
R^k_{\omega,\sigma_+,\sigma_+}\twoheadrightarrow R^k_{\omega,\omega,\sigma_+},
\qquad
R^k_{\omega,\sigma_-,\sigma_-}\twoheadrightarrow R^k_{\omega,\omega,\sigma_-},
\]
而要得到逆向系统，我们需要把 $R^k_{\omega,\omega,\sigma_-}$ 与 $R^k_{\omega,\omega,\sigma_+}$ 等同。然而由于奇点的存在，我们没有典范的等同。事实上存在两个完全自然的等同：一个记作 $\Psi$，用 $\mathcal{P}_\varphi$ 中沿从 $y_+$ 经 $x$ 到 $y_-$ 的路径（如图 9 所示）的平行移动；另一个记作 $\Psi'$，来自沿从 $y_+$ 经 $x'$ 到 $y_-$ 的类似路径。
由于我们已经通过 $x$ 把 $\Lambda_{y_+}$ 与 $\Lambda_{y_-}$ 等同，
\[
\Psi:R^k_{\omega,\omega,\sigma_+}\to R^k_{\omega,\omega,\sigma_-}
\]
由
\[
u\mapsto u,\qquad v\mapsto v,\qquad w\mapsto w
\]
给出。另一方面，由沿 $x'$ 的平行移动给出的 $\Lambda_{y_+}\to\Lambda_{y_-}$ 的映射恰好是 $T_\gamma$ 的逆，即 $e_1\mapsto e_1$、$e_2\mapsto e_2-e_1$，所以
\[
\Psi':R^k_{\omega,\omega,\sigma_+}\to R^k_{\omega,\omega,\sigma_-}
\]
由
\[
u\mapsto uw^{-1},\qquad v\mapsto vw,\qquad w\mapsto w
\]
给出。

% [原文 p.270]
现在在这两种情形中，我们都可以取逆向极限。此时逆向极限就是环的纤维积
\[
R^k_\omega:=R^k_{\omega,\sigma_+,\sigma_+}\times_{R^k_{\omega,\omega,\sigma_-}}R^k_{\omega,\sigma_-,\sigma_-},
\]
其中映射 $R^k_{\omega,\sigma_+,\sigma_+}\to R^k_{\omega,\omega,\sigma_-}$ 是自然满射 $R^k_{\omega,\sigma_+,\sigma_+}\twoheadrightarrow R^k_{\omega,\omega,\sigma_+}$ 与 $\Psi$ 或 $\Psi'$ 的复合，而映射 $R^k_{\omega,\sigma_-,\sigma_-}\to R^k_{\omega,\omega,\sigma_-}$ 就是自然满射。用 $\Psi$ 作纤维积时，我们得到
\[
R^k_\omega=\Bbbk[U,V,W^{\pm1},t]/(UV-t,t^{k+1}),
\]
其中
\[
U=(u,u),\qquad V=(v,v),\qquad W=(w,w),\qquad t=(uv,uv);
\]
而用 $\Psi'$ 作纤维积时，我们得到
\[
R^k_\omega=\Bbbk[U,V,W^{\pm1},t]/(UV-tW,t^{k+1}),
\]
其中
\[
U=(uw,u),\qquad V=(v,vw),\qquad W=(w,w),\qquad t=(uv,uv).
\]
没有自然的方式把这两个环等同，这是一个必须克服的根本问题。
解决办法是把环 $R^k_{\omega,\omega,\sigma_\pm}$ 在 $1+w^{-1}$（等价地，$1+w$）处局部化，并把映射 $\Psi$ 与 $\Psi'$ 各与 $(R^k_{\omega,\omega,\sigma_-})_{1+w^{-1}}$ 的一个自同构复合来修改它们。我们取新的 $\Psi$ 为旧的 $\Psi$ 与自同构
\begin{align*}
u&\mapsto u(1+w^{-1}),\\
v&\mapsto v(1+w^{-1})^{-1},\\
w&\mapsto w
\end{align*}
的复合，并取新的 $\Psi'$ 为旧的 $\Psi'$ 与自同构
\begin{align*}
u&\mapsto u(1+w),\\
v&\mapsto v(1+w)^{-1},\\
w&\mapsto w
\end{align*}
的复合。这些新的映射 $\Psi$、$\Psi'$ 由
\begin{align*}
\Psi&:u\mapsto u(1+w^{-1}),\quad v\mapsto v(1+w^{-1})^{-1},\quad w\mapsto w;\\
\Psi'&:u\mapsto uw^{-1}(1+w),\quad v\mapsto vw(1+w)^{-1},\quad w\mapsto w
\end{align*}
给出。注意这两个映射现在实际上重合了。于是纤维积现在是良定义的，并且可以看出（将在 §6.2.6 中证明），利用上述 $\Psi$ 或 $\Psi'$ 的选择，
\[
R^k_{\omega,\sigma_+,\sigma_+}\times_{(R^k_{\omega,\omega,\sigma_-})_{1+w^{-1}}}R^k_{\omega,\sigma_-,\sigma_-}
\]
现在是环
\[
R^k_\omega:=\Bbbk[U,V,W^{\pm1},t]/(UV+(1+W^{-1})t,t^{k+1}),
\]
其中
\[
U=(u,u(1+w^{-1})),\qquad V=(v(1+w^{-1}),v),\qquad W=(w,w),\qquad t=(uv,uv).
\]
注意 $\operatorname{Spec}\Bbbk[U,V,W^{\pm1},t]/(UV+(1+W^{-1})t)$ 现在定义了 $\operatorname{Spec}\Bbbk[U,V,W^{\pm1}]/(UV)$ 的一个光滑化，但该光滑化的全空间在 $U=V=t=0$、$W=-1$ 处是奇异的。这个奇点反映了仿射结构的奇点。

% [Tutor Note: p.270 的核心是：用 $\Psi$ 或 $\Psi'$ 作纤维积会得到两个不同的 $R^k_\omega$（分别为 $\Bbbk[U,V,W^{\pm1},t]/(UV-t,t^{k+1})$ 与 $(UV-tW,t^{k+1})$）；在 $1+w^{-1}$ 处局部化并用自同构 $u\mapsto u(1+w^{-1})$、$v\mapsto v(1+w^{-1})^{-1}$ 修正后两映射重合，最终得到 $R^k_\omega=\Bbbk[U,V,W^{\pm1},t]/(UV+(1+W^{-1})t,t^{k+1})$。核对锚点：全空间在 $U=V=t=0,\ W=-1$ 处奇异。]
% [原文 p.271]
这个修正初看可能相当特殊，但事实上它是根本性的。理解它的方式如下。想象从奇点出发沿边 $\omega$ 向相反方向延伸的两条射线。把图 9 中向右的射线标上 $1+w^{-1}$，向左的标上 $1+w$。这告诉我们：若需要越过其中一条射线来比较两个环，就用一个由所附函数定义的自同构来扭转这两个环之间的映射。
为了在 $\omega$ 之外也能一致地粘合，我们需要想象这些射线在 $B$ 中无限延伸。然而这些射线会彼此相交，破坏粘合的相容性，因为相应的粘合扭转不交换。如何修正这一问题正是 Kontsevich 与 Soibelman 在 [70] 中解决的关键问题。办法是从旧射线的交点出发添加带有新附函数的新射线，使相容性得以恢复。这些新射线又会产生新的交点，于是添加更多射线。最终对每个 $k$ 得到描述 $k$ 阶形变的有限射线集。当 $k\to\infty$ 时，我们通常预期射线的数目无界增长。
在接下来的几节中，我们将把整个过程精确化。

\subsection{结构}
本节我们将把上一小节所勾勒的“修正粘合”所需的数据形式化。指定粘合修正方式的数据称为结构（structure），它由一组射线组成。射线是一个几何对象（本质上是 $B$ 中的一条射线，我们称之为裸射线），连同一个函数，后者将指定越过该射线时的粘合。我们先定义裸射线。
本节我们固定 $(B,\mathcal{P},\varphi)$，其中 $B$ 是带奇点的整热带流形。此外我们假设 $B$ 是二维的、正的、简单的。我们允许 $B$ 有边界或非紧。令
\[
i:B_0\hookrightarrow B
\]
为包含，并令
\[
\Delta:=B\setminus B_0.
\]

\begin{definition}
$B$ 上的一个裸射线 $\mathfrak{d}$ 是一个映射
\[
\mathfrak{d}:I_{\mathfrak{d}}\to B,
\]
其中区间 $I_{\mathfrak{d}}$ 是 $[0,+\infty)$ 或 $[0,L_{\mathfrak{d}}]$（对某个 $L_{\mathfrak{d}}>0$），且满足以下性质：
\begin{enumerate}[label=(\arabic*),leftmargin=2.2em]
\item $\mathfrak{d}$ 是连续浸入，即在 $I_{\mathfrak{d}}$ 上局部单射。此外，它在 $\mathfrak{d}^{-1}(B_0)$ 上可微。
\item 若 $\mathfrak{d}(0)\in B_0$，则 $\mathfrak{d}(0)$ 在某个整仿射坐标卡下具有有理坐标。
\item 存在一个非零的整体截面
\[
\bar m_{\mathfrak{d}}\in\Gamma(I_{\mathfrak{d}},\mathfrak{d}^{-1}i_*\Lambda).
\]
特别地，若 $x\in\mathfrak{d}^{-1}(B_0)$，则 $\mathfrak{d}^{-1}i_*\Lambda$ 在 $x$ 处的茎与 $\Lambda$ 在 $\mathfrak{d}(x)$ 处的茎典范等同，于是这指定了 $\mathfrak{d}(x)$ 处的一个切向量 $\bar m_{\mathfrak{d},x}$。我们进而要求：若 $t$ 是 $I_{\mathfrak{d}}$ 上的坐标，则对任意点 $x\in\mathfrak{d}^{-1}(B_0)$，在 $x$ 处有
\[
\mathfrak{d}^*(\partial/\partial t)=-\bar m_{\mathfrak{d},x}.
\]
所以我们可以把 $-\bar m_{\mathfrak{d},x}$ 看作指定了速度。

% [原文 p.272]
\item 若 $I_{\mathfrak{d}}=[0,L_{\mathfrak{d}}]$，则 $\mathfrak{d}(L_{\mathfrak{d}})\in\partial B$ 且 $\mathfrak{d}([0,L_{\mathfrak{d}}))\subseteq B\setminus\partial B$。若 $I_{\mathfrak{d}}=[0,+\infty)$，则 $\mathfrak{d}(I_{\mathfrak{d}})\subseteq B\setminus\partial B$。于是特别地，若 $B$ 没有边界，则 $I_{\mathfrak{d}}=[0,+\infty)$。
\end{enumerate}
\end{definition}

让我们探讨这意味着什么。首先，条件 $\mathfrak{d}^*(\partial/\partial t)=-\bar m_{\mathfrak{d},x}$ 告诉我们：对任何满足 $\mathfrak{d}(U)\subseteq B_0$ 的连通开集 $U\subseteq I_{\mathfrak{d}}$，$U$ 的像是局部具有有理斜率的线段，因为它有局部常值且整的切向量。
其次，在 $\mathfrak{d}(x)\in\Delta$ 的点 $x\in I_{\mathfrak{d}}$ 处的行为受到限制。回忆在绕 $\Delta$ 中点 $p$ 的、以 $y$ 为基点的环路中，局部系统 $\Lambda$ 的单值性 $T$ 在茎 $\Lambda_y$ 的某组基 $e_1,e_2$ 下由 $T(e_1)=e_1$、$T(e_2)=e_1+e_2$ 给出。这意味着若 $V$ 是 $p$ 的小开邻域，则 $\Gamma(V\setminus\Delta,\Lambda)$ 可以认同为 $\Lambda_y$ 的单值不变元素，故
\[
\Gamma(V\setminus\Delta,\Lambda)\cong\Z e_1.
\]
于是按前推的定义 $\Gamma(V,i_*\Lambda)=\Z e_1$，特别地对 $x\in I_{\mathfrak{d}}$ 的小开邻域 $U$，有 $\Gamma(U,\mathfrak{d}^{-1}i_*\Lambda)=\Z e_1$。特别地，对任何 $x'\in U\setminus\{x\}$，$\bar m_{\mathfrak{d},x'}$ 必须属于 $\Lambda_{\mathfrak{d}(x')}$ 中在绕 $p$ 的环路的平行移动所给单值性下不变的元素构成的子群。
于是，当 $\mathfrak{d}$ 的像经过 $B$ 的奇点时，它必须沿唯一的单值不变切方向经过。但我们确切知道这些单值不变切方向是什么：它们由含该奇点的 $\mathcal{P}$ 的边给出。确实，我们在例 1.28 的 (4) 中明确看到了这一点。总结起来：$\mathfrak{d}$ 的像局部是具有有理斜率的直线，而当它经过一个奇点时，它与含该奇点的 $\mathcal{P}$ 的边重合。
注意构造裸射线也很容易。取一点 $y\in B_0\setminus\partial B$，它在 $y$ 附近的任一（等价地，所有）整仿射坐标卡中都有有理坐标。取 $\bar m\in\Lambda_y$，并把 $\mathfrak{d}:I_{\mathfrak{d}}\to B$ 定义为微分方程
\[
\mathfrak{d}'(x)=-\bar m_x
\]
的解，其中 $\bar m_x$ 表示 $\bar m$ 沿由 $\mathfrak{d}$ 决定的路径的平行移动。换句话说，我们只需把 $\mathfrak{d}$ 局部延拓为一条有理斜率的直线，一旦给定初始斜率 $-\bar m$ 它就完全确定。到达 $B$ 的边界时我们就停止。
有一个问题：若我们撞上 $B$ 的奇点 $p$ 会怎样？若 $\mathfrak{d}$ 在 $p$ 附近与经过 $p$ 的 $\mathcal{P}$ 的边重合，那就没问题：此时我们直接把 $\mathfrak{d}$ 延拓穿过奇点。然而我们不允许以任何其它方式穿过 $p$。为了避免这一点，我们作如下假设：

\begin{assumption}
对 $p\in\Delta$，设 $\omega\in\mathcal{P}$ 是含 $p$ 的边。注意 $\omega$ 被认同为 $\R$ 中的一个格线段，因此要求 $p$ 是 $\omega$ 中的非有理点是有意义的。我们假设所有这样的 $p$ 确实是非有理的。由于 $p$ 在 $\omega$ 中的选择在构造 1.25 中是任意的，我们总可以这样假设。
\end{assumption}

作为这一假设的推论，我们永远不必担心 $\mathfrak{d}$ 以非法方式经过奇点：由于 $\mathfrak{d}(0)$ 是有理的且 $\mathfrak{d}$ 具有有理斜率，只要 $\mathfrak{d}$ 与边 $\omega$ 横截，$\mathfrak{d}$ 的像总在有理点处与 $\omega$ 相交。

% [原文 p.273]
若我们想定义满足 $\mathfrak{d}(0)\in\Delta$ 的裸射线 $\mathfrak{d}$，则对 $\mathfrak{d}$ 的初始方向只有两个选择：它必须与含 $\mathfrak{d}(0)$ 的边相切。然后我们可以像之前那样构造这条射线。即使 $\mathfrak{d}(0)$ 不是有理点，由该假设它仍然只沿单值不变方向经过奇点。
现在我们可以定义射线：

\begin{definition}
$B$ 中的一条射线（ray）$(\mathfrak{d},f_{\mathfrak{d}})$ 由两部分数据组成：
\begin{itemize}[leftmargin=2em]
\item $\mathfrak{d}:I_{\mathfrak{d}}\to B$，一条裸射线；
\item $f_{\mathfrak{d}}=1+c_mz^m$，其中 $c_m\in\Bbbk$，$m$ 是 $\Gamma(I_{\mathfrak{d}},\mathfrak{d}^{-1}i_*\mathcal{P}_\varphi)$ 的一个截面，使得 $\bar m\in\Gamma(I_{\mathfrak{d}},\mathfrak{d}^{-1}i_*\Lambda)$（见定义 6.11）与 $\bar m_{\mathfrak{d}}$ 成正常数比例。此外，对每个满足 $\mathfrak{d}(x)\in B_0$ 的 $x\in I_{\mathfrak{d}}$，$m$ 在 $\mathcal{P}_{\varphi,\mathfrak{d}(x)}$ 中的芽 $m_x$ 属于幺半群 $\mathcal{P}_{\varphi,\mathfrak{d}(x)}$。
\end{itemize}
\end{definition}

% [Tutor Note: 本小节用弗拉克体字母：裸射线/射线记作 $\mathfrak{d}$（原书 \mathfrak{d}），$\operatorname{Chambers}(\mathcal{S},k)$ 的元素记作 $\mathfrak{u}$；结构中 $\mathcal{S}$ 为花体 S、$\mathcal{P}$ 为多面体分解。LaTeX 中分别用 \texttt{\textbackslash mathfrak\{d\}}、\texttt{\textbackslash mathfrak\{u\}}、\texttt{\textbackslash mathcal\{S\}}。]
让我们再一次考察这个定义。我们把一个多项式 $1+c_mz^m$ 附到一条裸射线上，其中 $m$ 是 $\Gamma(I_{\mathfrak{d}},\mathfrak{d}^{-1}i_*\mathcal{P}_\varphi)$ 的一个截面。现在若 $x\in I_{\mathfrak{d}}$ 满足 $\mathfrak{d}(x)\in\Delta$，设 $V$ 是 $\mathfrak{d}(x)$ 的小邻域。则由多值分片线性函数的定义，$\varphi$ 可以在 $V$ 上由单值函数代表。于是 $\mathcal{P}_\varphi$ 在 $V\cap B_0$ 上局部地分裂为 $\Z\oplus\Lambda$，从而 $i_*\mathcal{P}_\varphi$ 在 $V$ 上的一个截面可以写成 $(r,\bar m)$，其中 $\bar m\in\Gamma(V\cap B_0,\Lambda)$ 与经过奇点 $\mathfrak{d}(x)$ 的 $\mathcal{P}$ 的边相切。于是类似地，我们可以在 $x$ 的一个邻域 $U\subseteq I_{\mathfrak{d}}$ 中把 $m$ 写成 $(r,\bar m)$。特别地，$\bar m$ 总可以（至少在局部）提升为 $\mathfrak{d}^{-1}i_*\mathcal{P}_\varphi$ 的一个截面。
因此，截面 $m$ 可以简单地通过在某点 $x\in I_{\mathfrak{d}}$（满足 $\mathfrak{d}(x)\in B_0$）处给出 $m\in\mathcal{P}_{\varphi,\mathfrak{d}(x)}$ 来指定。我们要求 $\bar m$ 在 $\mathfrak{d}(x)$ 处与 $\bar m_{\mathfrak{d}}$ 成正比例。然后由平行移动，这个芽的选择就延拓为 $\mathfrak{d}^{-1}i_*\mathcal{P}_\varphi$ 的一个截面。
然而我们还有额外的要求：对所有的 $x\in\mathfrak{d}^{-1}(B_0)$ 都有 $m\in\mathcal{P}_{\varphi,\mathfrak{d}(x)}$。由于 $\mathfrak{d}$ 把 $I_{\mathfrak{d}}$ 映到一条可能以非常复杂的方式绕 $B$ 缠绕的射线或线段，这个条件看起来可能极强。但这里有一个非常重要的点：随着 $x\in I_{\mathfrak{d}}$ 增大，$m$ 关于 $\mathfrak{d}(x)$ 所在的各个极大胞腔的阶会增大。这是 $\mathfrak{d}(x)$ 沿 $-\bar m$ 方向移动这一事实的推论，并由下面的引理精确表述：

\begin{lemma}
\begin{enumerate}[label=(\arabic*),leftmargin=2.2em]
\item 设 $x\in B_0$，$m\in\mathcal{P}_{\varphi,x}$，且 $\bar m$ 指向某个胞腔 $\sigma\in\mathcal{P}_{\max}$ 的内部。若 $\operatorname{ord}_\sigma m\ge0$，则 $m\in\mathcal{P}_{\varphi,x}$。
\item 设 $x\in B_0\setminus\partial B_0$，$m\in\mathcal{P}_{\varphi,x}$。设 $\sigma_\pm\in\mathcal{P}_{\max}$ 是含 $x$ 的两个胞腔，具有性质：切向量 $-\bar m\in\Lambda_x$ 指向 $\sigma_-$ 内部，而 $\bar m$ 指向 $\sigma_+$ 内部。则对任何满足 $x\in\sigma$ 的 $\sigma\in\mathcal{P}_{\max}$，有
\[
\operatorname{ord}_\sigma(m)-\operatorname{ord}_{\sigma_+}(m)\ge0,\qquad
\operatorname{ord}_{\sigma_-}(m)-\operatorname{ord}_\sigma(m)\ge0.
\]
此外，若 $\bar m$ 不与 $\sigma\cap\sigma_+$ 相切，则
\[
\operatorname{ord}_\sigma(m)-\operatorname{ord}_{\sigma_+}(m)>0.
\]
\end{enumerate}
\end{lemma}

\begin{proof}
取 $x$ 的一个邻域 $U$ 与 $\varphi$ 在 $U$ 上的代表 $\varphi_U$，并对满足 $x\in\sigma$ 的 $\sigma\in\mathcal{P}_{\max}$ 令 $\varphi_{U,\sigma}\in\check\Lambda_x$ 为 $\varphi_U|_\sigma$ 的斜率。按惯例用 $\varphi_{U,x}(\bar m)=\varphi_{U,\sigma}(\bar m)$（若 $\bar m\in T_x\sigma$）定义 $\varphi_{U,x}:|\Sigma_x|\to\R$。对 (1)，注意 $\bar m$ 指向 $\sigma$ 内部蕴含 $\bar m\in|\Sigma_x|$，且 $\varphi_{U,x}(\bar m)=\varphi_{U,\sigma}(\bar m)$。于是 $\operatorname{ord}_\sigma(m)\ge0$ 蕴含（用 $\varphi_U$ 给的分裂写成 $m=(r,\bar m)$）$r\ge\varphi_{U,\sigma}(\bar m)=\varphi_{U,x}(\bar m)$。所以 $m\in\mathcal{P}_{\varphi,x}$。
对 (2)，注意由于 $x\notin\partial B$，$\Sigma_x$ 是完全扇。此外
\[
\operatorname{ord}_\sigma(r,\bar m)=r-\varphi_{U,\sigma}(\bar m).
\]
由于 $\bar m$ 指向 $\sigma_+$ 内部，由 $\varphi_{U,x}$ 的严格凸性，对所有含 $x$ 的 $\sigma\in\mathcal{P}_{\max}$ 有 $\varphi_{U,\sigma}(\bar m)\le\varphi_{U,\sigma_+}(\bar m)$。于是
\[
\operatorname{ord}_\sigma(m)\ge\operatorname{ord}_{\sigma_+}(m).
\]
对 $-\bar m$ 应用同样的论证给出 $\varphi_{U,\sigma}(-\bar m)\le\varphi_{U,\sigma_-}(-\bar m)$，即
\[
\operatorname{ord}_\sigma(m)\le\operatorname{ord}_{\sigma_-}(m).
\]
最后，若 $\bar m$ 不与 $\sigma\cap\sigma_+$ 相切，则由 $\varphi_{U,x}$ 的严格凸性有 $\varphi_{U,\sigma}(\bar m)\neq\varphi_{U,\sigma_+}(\bar m)$，故得最后一条。
\end{proof}

% [原文 p.274]
\begin{definition}
设 $\mathfrak{d}:I_{\mathfrak{d}}\to B$ 是一条裸射线，并设 $m\in\Gamma(I_{\mathfrak{d}},\mathfrak{d}^{-1}i_*\mathcal{P}_\varphi)$ 是一个截面，使得 $\bar m$ 是 $\bar m_{\mathfrak{d}}$ 的正常数倍。则定义 $\operatorname{ord}_m:\mathfrak{d}^{-1}(B_0)\to\Z$ 为
\[
\operatorname{ord}_m(x)=\sup\{\operatorname{ord}_\sigma(m_x)\mid \mathfrak{d}(x)\in\sigma\in\mathcal{P}_{\max}\}.
\]
引理 6.19 蕴含这是一个阶梯函数，每当 $\mathfrak{d}$ 进入一个新的极大胞腔时它就增大。
\end{definition}

我们现在看到，指定射线事实上很容易：

\begin{proposition}
设 $\mathfrak{d}:I_{\mathfrak{d}}\to B$ 是一条裸射线。若 $\mathfrak{d}(0)\in B_0$，设 $m\in\mathcal{P}_{\varphi,\mathfrak{d}(0)}$ 使得 $\bar m$ 是 $\mathfrak{d}'(0)\in\Lambda_{\mathfrak{d}(0)}$ 的负常数倍。若 $\mathfrak{d}(0)\in\Delta$，取 $\epsilon>0$ 足够小，使 $\mathfrak{d}([0,\epsilon])$ 含于 $\mathcal{P}$ 的某条边的内部，并设 $m\in\mathcal{P}_{\varphi,\mathfrak{d}(\epsilon)}$ 使得 $\bar m$ 是 $\mathfrak{d}'(\epsilon)\in\Lambda_{\mathfrak{d}(\epsilon)}$ 的负常数倍。则：
\begin{enumerate}[label=(\arabic*),leftmargin=2.2em]
\item $m$ 延拓为 $\Gamma(I_{\mathfrak{d}},\mathfrak{d}^{-1}i_*\mathcal{P}_\varphi)$ 的一个截面 $m$，使得对所有的 $x\in\mathfrak{d}^{-1}(B_0)$，$m$ 定义一个芽 $m_x\in\mathcal{P}_{\varphi,\mathfrak{d}(x)}\subseteq\mathcal{P}_{\varphi,\mathfrak{d}(x)}$。
\item 对每个非负整数 $k$，存在 $N\ge0$，使得若 $x\in[N,+\infty)\cap\mathfrak{d}^{-1}(B_0)$，$\tau\in\mathcal{P}$ 是唯一使 $\mathfrak{d}(x)\in\operatorname{Int}(\tau)$ 的胞腔，$\sigma\in\mathcal{P}_{\max}$ 是含 $\tau$ 的胞腔，则 $m_x\in I^k_{\tau,\tau,\sigma}$。
\end{enumerate}
\end{proposition}

\begin{proof}
芽 $m$ 经由平行移动延拓为 $\Gamma(I_{\mathfrak{d}},\mathfrak{d}^{-1}i_*\mathcal{P}_\varphi)$ 的一个截面；定义 6.18 之后的讨论表明，这个截面能跨越 $\mathfrak{d}^{-1}(\Delta)$ 的点，正是因为 $\mathfrak{d}$ 本身已经是一条裸射线。
$m_0\in\mathcal{P}_{\varphi,\mathfrak{d}(0)}$（或当 $\mathfrak{d}(0)\in\Delta$ 时 $m_\epsilon\in\mathcal{P}_{\varphi,\mathfrak{d}(\epsilon)}$）这一说法蕴含对含 $\mathfrak{d}(0)$ 的所有 $\sigma\in\mathcal{P}_{\max}$ 有 $\operatorname{ord}_\sigma m_0\ge0$（或对含 $\mathfrak{d}(\epsilon)$ 的所有 $\sigma\in\mathcal{P}_{\max}$ 有 $\operatorname{ord}_\sigma m_\epsilon\ge0$）。但引理 6.19 的 (2) 接着告诉我们这一点会继续保持，即对所有的 $x\in\mathfrak{d}^{-1}(B_0)$，在 $\mathfrak{d}(x)\in\sigma\in\mathcal{P}_{\max}$ 时都有 $\operatorname{ord}_\sigma(m_x)\ge0$。于是由同一条引理的 (1)，对所有的 $x\in\mathfrak{d}^{-1}(B_0)$ 有 $m_x\in\mathcal{P}_{\varphi,\mathfrak{d}(x)}$。
对 (2)，我们只需选取 $N$ 使得对所有的 $x\ge N$ 有 $\operatorname{ord}_m(x)>k$。
\end{proof}

现在我们终于可以定义结构的概念了。

\begin{definition}
一个结构 $\mathcal{S}$ 是一族射线，满足以下性质：
\begin{enumerate}[label=(\arabic*),leftmargin=2.2em]
\item 对每个整数 $k\ge0$，令 $\mathcal{S}[k]$ 为 $\mathcal{S}$ 中满足如下条件的射线 $(\mathfrak{d},1+c_mz^m)$ 的集合：存在 $x\in\mathfrak{d}^{-1}(B_0)$ 使 $\operatorname{ord}_m(x)\le k$。则我们要求对每个 $k$，$\mathcal{S}[k]$ 是有限的。

% [原文 p.275]
\begin{figure}[htbp]
\centering
\includegraphics[width=0.46\textwidth]{f6_10.png}
\figtitle{10}{}
\end{figure}
\item 若 $\omega\in\mathcal{P}$ 是含奇点 $p$ 的边，则 $\mathcal{S}$ 含两条射线 $\mathfrak{d}_{p,\pm}$，使得 $\mathfrak{d}_{p,\pm}(0)=p$，且对接近零的 $\epsilon$ 有 $\mathfrak{d}'_{p,\pm}(\epsilon)=\pm\bar m$，其中 $\bar m$ 是与 $\omega$ 相切的本原切向量。最后，
\[
f_{\mathfrak{d}_{p,\pm}}=1+z^{m_\mp},
\]
其中 $m_\mp\in\mathcal{P}_{\varphi,x}$（对任意 $x\in\operatorname{Int}(\omega)\setminus\{p\}$）由条件 $\bar m_\pm=\pm\bar m$ 与“对任何含 $\omega$ 的极大胞腔 $\sigma$ 有 $\operatorname{ord}_\sigma m_\pm=0$”唯一确定。这样的射线称为初始射线。我们把 $\mathcal{S}$ 的初始射线集记作 $\operatorname{Initial}(\mathcal{S})$。对 $\mathfrak{d}\in\operatorname{Initial}(\mathcal{S})$ 且 $\mathfrak{d}(0)\in\omega$ 为边的情形，我们定义 $I^{\mathrm{init}}_{\mathfrak{d}}$ 为最大的子区间 $[0,L]\subseteq I_{\mathfrak{d}}$，使得 $\mathfrak{d}([0,L])\subseteq\omega$。
\item 若 $(\mathfrak{d},f_{\mathfrak{d}})\in\mathcal{S}$ 不是初始射线，且 $f_{\mathfrak{d}}=1+c_{\mathfrak{d}}z^{m_{\mathfrak{d}}}$，则对每个 $x\in\mathfrak{d}^{-1}(B_0)$ 有 $\operatorname{ord}_{m_{\mathfrak{d}}}(x)>0$。
\end{enumerate}
\end{definition}

\begin{example}
考虑图 10 中所描绘的整热带流形。图中上方图形描绘了把 $B$ 减去图中切口后嵌入 $\R^2$ 的嵌入，它描述了切口之外的仿射结构。此外，在顶点 $(-1,0)$ 与 $(0,-1)$ 处的扇结构如图中所示，它们指定了跨越这些切口的仿射结构。容易验证 $B$ 是正的且简单的。取函数 $\varphi$ 在 $\sigma_1$ 上斜率为 $(2,1)$、在 $\sigma_2$ 上为 $(1,2)$、在 $\sigma_3$ 上为 $(0,2)$、在 $\sigma_4$ 上为 $(0,0)$、在 $\sigma_5$ 上为 $(2,0)$，并取其在原点处值为零；容易验证这是严格凸的、良定义的、且跨越切口连续。图 11 现在描绘了一个结构 $\mathcal{S}$，它由五条有限长度的射线组成。把 $\Lambda$ 的元素写成 $\Z^2$ 的元素，使用图 10 中所描绘的主坐标卡，并使用由上述 $\varphi$ 的选择给出的 $\mathcal{P}_\varphi=\Z\oplus\Lambda$ 的分裂，

% [原文 p.276]
\begin{figure}[htbp]
\centering
\includegraphics[width=0.46\textwidth]{f6_11.png}
\figtitle{11}{}
\end{figure}
我们有从每个奇点出发的两条射线（即初始射线），以及一条额外的射线，它从原点出发沿东北方向前进。
\end{example}

% [Tutor Note: 图 10 与图 11 均无图注文字（原文仅编号），按 Lead 指示用 $\figtitle{10}{}$、$\figtitle{11}{}$；图 10 裁自 PDF p.290（含上方的扇结构与下方的两块阴影），图 11 裁自 PDF p.291（五条射线及其附函数 $1+z^{(\cdot)}$ 的标号）。]
\subsection{相容结构与粘合}
我们把一个结构看作给出了一组用于粘合诸环 $R^k_{\tau_1,\tau_2,\sigma}$ 的拷贝的数据。我们现在说明这如何进行，然后定义保证这种粘合有意义的必要条件。
对这项构造，我们固定一个结构 $\mathcal{S}$ 与整数 $k\ge0$。虽然 $\mathcal{S}$ 可能含无限多条射线，但 $\mathcal{S}[k]$ 是有限的。此外，对每个 $\mathfrak{d}\in\mathcal{S}[k]$，设 $N^k_{\mathfrak{d}}>0$ 是命题 6.21 的 (2) 所保证的数。

\begin{definition}
设 $\mathcal{P}_k$ 是 $\mathcal{P}$ 的一个加细，满足以下性质：
(1) $\mathcal{P}_k$ 的元素是具有理顶点的凸多面体。
(2) 对 $\mathfrak{d}\in\mathcal{S}[k]$，$\mathfrak{d}([0,N^k_{\mathfrak{d}}])$ 含于 $\mathcal{P}_k$ 的若干边的并，而 $\mathfrak{d}(N^k_{\mathfrak{d}})$ 是 $\mathcal{P}_k$ 的顶点。
(3) 对 $\mathfrak{d}\in\mathcal{S}[k]\setminus\operatorname{Initial}(\mathcal{S})$，$\mathfrak{d}(0)$ 是 $\mathcal{P}_k$ 的顶点。
\end{definition}

\begin{example}
在例 6.23 中，我们可以对所有的 $k$ 取 $\mathcal{P}_k=\mathcal{P}$。
\end{example}

选定了 $\mathcal{P}_k$ 之后（我们的构造不依赖于具体选择），定义 $\operatorname{Chambers}(\mathcal{S},k)$ 为 $\mathcal{P}_k$ 的极大胞腔的集合。对每个 $\mathfrak{u}\in\operatorname{Chambers}(\mathcal{S},k)$，存在唯一的 $\sigma\in\mathcal{P}_{\max}$ 包含 $\mathfrak{u}$，我们把它写作 $\sigma_{\mathfrak{u}}$。
我们现在如下定义一个范畴 $\operatorname{Glue}(\mathcal{S},k)$：
\begin{itemize}[leftmargin=2em]
\item $\operatorname{Glue}(\mathcal{S},k)$ 的对象是三元组 $(\omega,\tau,\mathfrak{u})$，其中
\[
\omega,\tau\in\mathcal{P},\quad \mathfrak{u}\in\operatorname{Chambers}(\mathcal{S},k),\quad \text{且}\quad \omega\subseteq\tau,\ \ \omega\cap\mathfrak{u}\neq\varnothing,\ \ \tau\subseteq\sigma_{\mathfrak{u}}.
\]
\end{itemize}
% [Tutor Note: 本块（c06b）末尾的 itemize 仅含“Glue(𝒮,k) 的对象”这一条，已在本块闭合；c06d 续写时请重新 \begin{itemize} 写 p.277 首条“至多一个态射……”（原文用 \underline{Glue} 记范畴名，本译作 \texttt{\textbackslash operatorname\{Glue\}}）。另：本块 16 个定理类环境自动编号应为 6.10–6.25，需 c06a 恰含 9 个环境（6.1–6.9）；Assumption 6.17 依赖 preamble 新增的 assumption 环境。]
